
\minitoc 

% ========================================================================================================
% Introduction 

La simulation peut se définir comme la prédiction de l'évolution d'un système en résolvant des équations 
venues de domaines appliqués. 
La notion d'aléatoire vient du fait que l'on simule des phénomènes dont nous n'aurons que des informations imparfaites 
tels que des phénomènes de désintégration nucléaire ou de prévention des risques. Les utilisations 
possibles de la simulation sont : 
\begin{itemize}
    \item Étude d'un phénomène (ex : la trajectoire d'un avion dans des turbulences)
    \item L'estimation de caractéristiques d'un système (ex : déterminer les caractéristiques d'une simulation météo)
    \item Tester des hypothèses (ex : en médecine, on peut utiliser la simulation aléatoire pour déterminer à l'avance 
    les effets d'un médicament sur un patient.)
    \item La prise de décision dans un environement incertain. 
    \item Réaliser des expériences coûteuses voire impossibles (ex : effectuer des essais nucléaires informatiquement 
    en fonction des essais passés ou simuler une propagation virale dans une population). 
\end{itemize}
La simulation aléatoire est donc un outil rapide économique et contrôlé. 

\begin{definition}[Simuler une variable aléatoire]
    Simuler une variable aléatoire de loi $\mu$ consiste à produire un algorithme dont les réalisations donnent 
    des copies indépendantes d'une variable de loi $ \mu$. 
\end{definition}

La simulation aléatoire est impossible sans l'accès à des \emph{données aléatoires}. En pratique, on se contentera 
de données \emph{pseudo-aléatoires} tel que le nombre de secondes écoulées depuis le 01/01/1070, la température 
d'un processeur ou les dernières touches entrées avant l'exécution du programme. 
Ceci est obtenu en \texttt{Python} avec la fonction \texttt{rand()}.  Dans ce cours, on considère que chaque appel 
à cette fonction nous renvoie une variable iid de loi Uniforme sur $[0,1]$. 

Simuler une variable aléatoire consiste donc à construire un algorithme qui transforme un nombre fini de variables 
aléatoires idd de loi $ \mathcal{U}([0,1])$ en une variable aléatoire de loi $ \mu$ : 
    \[ \mathcal{U}([0,1]) \overset{\text{simulation}}{\longrightarrow} \mu \] 

% ========================================================================================================
% Simulation de variables aléatoires discrètes

\section{Simulation de variables aléatoires discrètes}

\begin{example}[Simuler une variable aléatoire de loi $ \mathcal{B}(p)$]
    Soit $X \sim \mathcal{U}([0,1])$ avec $p \in [0,1]$. Alors la variable $Y = \mathds{1}_{X \leqslant p}$  
    suit un loi de Bernouilli de paramètre $p$. 
        \[ \myP(Y = 1) = \myP(\mathds{1}_{X \leqslant p} = 1) = \myP(X \leqslant p) = p \]  
        \[ \myP(Y = 0) = 1 - p \] 
    On peut donc développer l'algorithme suivant :
    \begin{lstlisting}[language=Python]
def bernouilli(p): 
    x = rand()
    if u <= p: 
        return 1 
    return 0
    \end{lstlisting}
\end{example}

\begin{example}[Simuler une loi binomiale $ \mathcal{B}(n,p)$]
    Soient $(X_1, \dots, X_n)$ une suite de variables aléatoires iid de loi $ \mathcal{B}(n,p)$. 
    La variable $ Y := \sum_{n=1}^{n} X_i \sim \mathcal{B}(n,p)$ telle que : 
        \[ \myP(Y = k) = \binom{n}{k} p^k (1-p)^{n-k} \] 
    Il nous faut donc seulement générer la somme de $n$ variables aléatoires suivant une loi de Bernouilli. On a donc : 
        \begin{lstlisting}[language=Python]
def binom(n,p):
    s = 0 
    for i in range(n):
        s += 1 * (rand() <= p) 
    return s
    \end{lstlisting} 
\end{example}

On cherchera à développer des programmes qui terminent en temps fini. On préferera deux qui terminent 
en temps borné (comme \texttt{binom(n,p)}). 

\begin{proposition}
    Soit $(p_n)_{n \geqslant 1}$ une suite réelle positive telle que : 
        \[ \sum_{k=1}^{\infty} p_k = 1 \] 
    et $(x_n)$ une suite d'éléments d'un ensemble $E$. On définit la loi $ \mu$ par : 
        \[ \mu = \sum_{k=1}^{\infty} p_k \delta_{x_k} \] 
    si une variable aléatoire $X$ est de loi $ \mu$ on aura alors $ \myP(X = x_k) = p_k$. 
    On pose $f_0 = 0$ et $s_k = p_1 + \dots + p_k$. Soit $U$ une variable aléatoire de loi $ \mathcal{U}([0,1])$. 
    Alors la variable : 
        \[ Y = \sum_{k=1}^{\infty} x_n \times \mathds{1}_{s_{n-1} \leqslant U \leqslant s_k} \] 
    suit la loi $ \mu$. L'algorithme suivant simule une variable aléatoire de loi $ \mu$. 
    \begin{lstlisting}[language=Python]
def simulation_mu([p1, ..., pk],[x1, .., xk]): 
    u = rand() 
    s = p1 
    while u > s :
        i++ 
        s += pi 
    return xi 
    \end{lstlisting}  
\end{proposition}

\begin{proof}
    Pour tout $n \in \N$, on a :
        \[ \myP(s_{n-1} \leqslant U \leqslant s_n) = s_n - s_{n-1} = p_n \] 
    donc $ \myP(Y = x_n) = \myP(s_{n-1} \leqslant U \leqslant s_n) = p_n $.   
\end{proof}

\begin{example}[Simuler une loi de Poisson $ \mathcal{P}(\lambda)$ ]
    Pour rappel, si $X \sim \mathcal{P}(\lambda)$ alors : 
        \[ \forall k \in \N, \quad \myP(X = k) = e^{-\lambda} \frac{\lambda^k}{k !} \]  
    Posons $ p_k := e^{- \lambda} \frac{ \lambda^k}{k!}$. On remarque que : 
        \[ p_{k+1} = e^{- \lambda} \frac{ \lambda^{k+1}}{(k+1)!} = p_k \frac{ \lambda}{k+1} \] 
    On en déduis donc l'algorithme suivant : 

    \begin{lstlisting}[language=Python]
def poisson(lambda): 
    x = 0 
    p = exp(-\lamdba) 
    s = p 
    u = rand() 
    while u > s : 
        x += 1
        p = p * lambda / x 
        s += p 
    return x  
    \end{lstlisting}
\end{example}

\begin{example}[Loi uniforme]
    Simulons une loi uniforme sur $\{1, \dots, n\}$. 
    On a donc l'algorithme suivant : 

    \begin{lstlisting}[language=Python]
def uniforme(n):
    x = 1 
    s = x/n 
    u = rand() 
    while u > s : 
        x += 1
        s += 1/n 
    return x 

def uniforme(n):
    return ceil(n * rand(u))
    \end{lstlisting}
\end{example}

\begin{proposition}[Algorithme du Rejet]
    Soit $(E, \mathcal{E}, \mu)$ un espace probabilisé et $A \subset E$ un évènement tel que $\mu(A) > 0$. 
    La loi $\mu$ conditionnellement à $A$, notée $\mu_A$, est définie par :
    \[
        \mu_A(B) = \frac{\mu(A \cap B)}{\mu(A)}, \qquad B \in \mathcal{E}.
    \]
    Soit $(X_n)_{n \geqslant 1}$ une suite de variables aléatoires i.i.d. de loi $\mu$. 
    On pose :
    \[
        T := \inf \{ n \geqslant 1 \mid X_n \in A \},
    \]
    c'est-à-dire le rang du premier tirage appartenant à $A$.
    Alors :
    \begin{enumerate}
        \item $T \sim \mathcal{G}(\mu(A))$, c'est-à-dire
        \[
            \mathbb{P}(T = k) = (1 - \mu(A))^{k-1} \mu(A), \qquad k = 1,2,\dots
        \]
        \item $X_T \sim \mu_A$ (loi conditionnelle sur $A$)
        \item $X_T$ et $T$ sont indépendantes.
    \end{enumerate}
\end{proposition}

\begin{proof}
    Soit $B \in \mathcal{E}$ et $k \geqslant 1$. On calcule :
    \begin{align*}
        \mathbb{P}(X_T \in B, T = k) 
        &= \mathbb{P}(X_k \in B, T = k) \\
        &= \mathbb{P}(X_1 \notin A, \dots, X_{k-1} \notin A, X_k \in A, X_k \in B) \\
        &= \mathbb{P}(X_1 \notin A)^{k-1} \times \mathbb{P}(X_1 \in A \cap B) \\
        &= (1 - \mu(A))^{k-1} \mu(A) \frac{\mu(A \cap B)}{\mu(A)} \\
        &= \underbrace{(1 - \mu(A))^{k-1} \mu(A)}_{\mathbb{P}(T = k)}
           \underbrace{\frac{\mu(A \cap B)}{\mu(A)}}_{\mu_A(B)}.
    \end{align*}
    On en déduit :
    \begin{itemize}
        \item $\mathbb{P}(T = k) = (1 - \mu(A))^{k-1} \mu(A)$ (loi géométrique)
        \item $\mathbb{P}(X_T \in B) = \mu_A(B)$
        \item $\mathbb{P}(X_T \in B, T = k) = \mathbb{P}(T = k) \mathbb{P}(X_T \in B)$, donc $T$ et $X_T$ sont indépendantes.
    \end{itemize}
\end{proof}


\paragraph{Interprétation.}
L'algorithme du rejet décrit la procédure suivante :
\begin{enumerate}
    \item Tirer $X_1, X_2, \dots$ de manière indépendante selon la loi $\mu$.
    \item Rejeter tout tirage $X_n \notin A$ et recommencer.
    \item Accepter le premier $X_T \in A$ : sa loi est exactement $\mu_A$.
\end{enumerate}

Ainsi :
\begin{itemize}
    \item $T$ représente le nombre de tirages nécessaires avant d'obtenir un succès.
    \item $X_T$ est la valeur de ce premier succès, distribuée comme $\mu$ conditionnellement à $A$.
    \item $T$ et $X_T$ sont indépendants : le temps d'attente n'influence pas la valeur obtenue.
\end{itemize}

\begin{example}
    Si $E$ est l'ensemble des notes d'étudiants, $\mu$ la distribution des notes et $A = \{\text{note} > 15\}$, 
    alors :
    \begin{itemize}
        \item $T$ = nombre d'étudiants observés avant d'en trouver un avec une note $> 15$,
        \item $X_T$ = note de ce premier étudiant (loi conditionnelle sur $A$).
    \end{itemize}
    La loi de $T$ est géométrique et la loi de $X_T$ est la loi des notes sachant qu'elles sont $ > 15$.
\end{example}

\begin{example}
    On veut simuler une variable aléatoire $U$ de loi uniforme sur les nombres premiers inférieurs à $n$. 
    D'après la proposition précédente, on peut donc utiliser l'algorithme suivant : 
    \begin{lstlisting}[language=Python]
def unif_prem(n):
    x = uniforme(n)
    while not premier(x): 
        x = uniforme(n) 
    return x
    \end{lstlisting}
    Ainsi le premier nombre premier sur lequel on tombe est uniformément distribué sur l'ensemble des nombres premiers 
    inférieurs ou égaux à $n$. 
\end{example}


